\section{بار کاری ورودی/خروجی (\lr{fileio})}
\fig{fig/io.png}{توان عملیاتی نوشتن، تأخیر میانگین و \lr{P99} در آزمون \lr{fileio} (محور عمودی لگاریتمی)}
\begin{center}
\begin{tabular}{lrrr}
\toprule
محیط & \lr{Write (MiB/s)} & \lr{avg (ms)} & \lr{P99 (ms)} \\
\midrule
\lr{Bare Metal (WSL)} & \lr{9.26} & \lr{0.74} & \lr{1.96} \\
\lr{VM} & \lr{609.72} & \lr{0.01} & \lr{0.05} \\
\lr{Docker} & \lr{45.22} & \lr{0.15} & \lr{0.39} \\
\bottomrule
\end{tabular}
\end{center}
\textbf{تحلیل.} این بخش بیشترین تفاوت را نشان می‌دهد و ترتیب آن (\lr{VM} سریع‌تر از \lr{Docker} و \lr{Docker} سریع‌تر از \lr{Bare Metal}) خلاف انتظار «مجازی‌سازی = سربار» است. این نتیجه نباید به‌عنوان سربار خالص مجازی‌سازی تفسیر شود، بلکه نشان‌دهندهٔ تفاوت مسیر ذخیره‌سازی در هر محیط است:
\begin{itemize}
  \item \textbf{\lr{VM}:} تأخیر ۰٫۰۱ میلی‌ثانیه و توان ۶۰۰ \lr{MiB/s} بسیار بالاتر از آن است که یک نوشتن واقعی روی دیسک بتواند پاسخ دهد. به‌احتمال زیاد نوشتن‌ها در حافظهٔ نهان (\lr{cache}) سیستم‌عامل میهمان یا \lr{hypervisor}/میزبان جذب شده‌اند و به دیسک فیزیکی نرسیده‌اند. پس عدد \lr{VM} بیشتر سرعت \lr{cache} را نشان می‌دهد تا سرعت دیسک.
  \item \textbf{\lr{Docker}:} فایل‌ها در سیستم‌فایل لایه‌ای کانتینر (مثلاً \lr{overlayfs}) نوشته می‌شوند که یک لایهٔ اضافه دارد، اما از \lr{cache} هستهٔ میزبان استفاده می‌کند. نتیجه میانی است.
  \item \textbf{\lr{Bare Metal (WSL)}:} کندترین است، که با نوشتن روی مسیر ویندوز (مثلاً \code{/mnt/c/} از طریق پروتکل \lr{9P}) سازگار است، چون هر عملیات باید از مرز لینوکس به ویندوز عبور کند. این مورد را با بررسی مسیر محل اجرای \lr{sysbench} باید تأیید کرد.
\end{itemize}
پراکندگی بین تکرارها در \lr{fileio} کوچک است (کمتر از حدود ۷٪ در توان عملیاتی)، پس این تفاوت‌ها پایدارند، اما علت آن‌ها تفاوت سیستم ذخیره‌سازی و \lr{caching} است نه خود فناوری کانتینر یا \lr{VM}.
